% ===================================================================
\section{Method: Generalized Pair Specification}
\label{sec:method}

The Steersman's contrastive loss, as defined in
Paper~3~\cite{bielec2026bridge}, encodes the rhombic dodecahedron's
face-pair geometry through a fixed partition of channel indices into
co-planar and cross-planar sets.  We generalize this mechanism by
replacing the hard-coded RD partition with a configurable pair
specification interface.  The Steersman itself---its three control laws
(connectivity trending, directionality stagnation, stability
monitoring), spectral regularization target, and feedback dynamics---is
\emph{unchanged}.  The only modification is the pair specification fed
to the contrastive loss.

\subsection{The Pair Specification Interface}
\label{sec:pairs}

The function \texttt{\_compute\_pair\_indices}$(n)$ returns two
lists: \emph{co-axial} pairs $\mathcal{C}$ (channels that should
couple strongly) and \emph{cross-axial} pairs $\mathcal{X}$ (all
remaining distinct-index pairs).  The contrastive loss attracts
co-axial pairs and repels cross-axial pairs, operating on the
partition $(\mathcal{C}, \mathcal{X})$ identically to Paper~3.
The spectral loss maintains overall algebraic connectivity by
targeting a Fiedler eigenvalue $\lambda_2 \approx 0.1$.

Four geometric specifications derived from polytope vertex-opposition
structure are tested:

\paragraph{Octahedron ($n{=}4$).}
Two co-axial pairs encoding a 2-axis partition of the octahedron's
three vertex-opposition axes:
\begin{equation}
  \texttt{pairs}_{\text{oct}}(4) = \{(0,1),\; (2,3)\}
  \label{eq:oct_pairs}
\end{equation}
with $|\mathcal{C}| = 2$ and $|\mathcal{X}| = 4$.  This is the
minimum-channel-count geometry: only two co-axial pairs and the
smallest cross/co ratio ($4/2 = 2$), providing the strongest co-axial
signal per channel pair.

\paragraph{Rhombic dodecahedron ($n{=}6$).}
Three co-planar pairs from the RD's 3-axis face-pair geometry:
\begin{equation}
  \texttt{pairs}_{\text{RD}}(6) = \{(0,5),\; (1,4),\; (2,3)\}
  \label{eq:rd_pairs}
\end{equation}
with $|\mathcal{C}| = 3$ and $|\mathcal{X}| = 12$.  Each pair shares
a vanishing coordinate in the face-normal vectors, producing three
independent $2 \times 2$ blocks when compiled.

\paragraph{Tesseract ($n{=}8$).}
Four co-axial pairs aligned with the four coordinate axes of the
4D hypercube:
\begin{equation}
  \texttt{pairs}_{\text{tess}}(8) = \{(0,1),\; (2,3),\; (4,5),\;
    (6,7)\}
  \label{eq:tess_pairs}
\end{equation}
with $|\mathcal{C}| = 4$ and $|\mathcal{X}| = 24$.  If the
Steersman generalizes beyond 3D, this specification should produce
four independent $2 \times 2$ blocks---a $4{+}4$ eigenvalue split.

\paragraph{24-cell ($n{=}24$, D4 root polytope).}
Twelve co-axial pairs from the 24 vertices expressed as permutations
of $(\pm 1, \pm 1, 0, 0)$ in $\mathbb{R}^4$.  Vertices are indexed
by coordinate plane: $wx$ (indices 0--3), $wy$ (4--7), $wz$ (8--11),
$xy$ (12--15), $xz$ (16--19), $yz$ (20--23), with two antipodal
pairs per plane:
\begin{equation}
  \texttt{pairs}_{\text{24c}}(24) =
    \bigl\{(0,3),(1,2),\; (4,7),(5,6),\; (8,11),(9,10),\;
    (12,15),(13,14),\; (16,19),(17,18),\; (20,23),(21,22)\bigr\}
  \label{eq:24c_pairs}
\end{equation}
with $|\mathcal{C}| = 12$ and $|\mathcal{X}| = \binom{24}{2} - 12
= 264$.  The 24-cell is the densest 4D sphere-packing polytope
(the D4 root system), standing in the same relation to 4D as the
rhombic dodecahedron does to 3D.  If the Steersman scales, this
specification should produce twelve independent $2 \times 2$
blocks---a $12{+}12$ eigenvalue split.

The interface imposes no constraint on which indices are paired.
Any partition of $n$~channels into $\lfloor n/2 \rfloor$ pairs
constitutes a valid input.  This generality enables the non-geometric
controls described next.

\subsection{Non-Geometric Controls}
\label{sec:controls}

Two negative controls test whether the Steersman requires geometric
coherence in the pair specification or compiles arbitrary partitions.

\paragraph{Wrong-Labels (WL-001).}
Six channels partitioned into three pairs by random permutation
(seed~42), producing co-axial pairs that do not correspond to any
polytope's antipodal structure.  The Steersman receives the same
contrastive loss function with the same hyperparameters as the
standard $n{=}6$ RD experiment; only the pair specification differs.
If the Steersman is a brute-force topology programmer, random
pairs should produce block-diagonal structure.  If geometric
coherence is required, the random partition should fail.

\paragraph{Resonance (R-001).}
Six channels with co-axial pairs derived from number-theoretic
relationships among corpus primes:  Sophie~Germain
($2(11){+}1 {=} 23 \to$ channels $(0,2)$), consecutive primes
($29, 31 \to$ channels $(3,4)$), and shared residue class
($11 \equiv 17 \equiv 5 \pmod{6} \to$ channels $(0,5)$).  These
pairs carry genuine algebraic structure---prime-theoretic bonds---but
no geometric embedding in any spatial dimension.  This control tests
whether mathematical structure in general, or spatial structure
specifically, is the operative requirement.

\subsection{Emanation Architecture (E-001)}
\label{sec:emanation_method}

The emanation architecture introduces hierarchical bridge structure as
a variant of the standard per-layer independent bridges.  A single
master bridge matrix $\mathbf{M}_{\text{master}}$ is shared across all
transformer layers, with per-layer offset matrices capturing local
variation:
\begin{equation}
  \mathbf{M}^{(l)} = \mathbf{M}_{\text{master}} +
    \boldsymbol{\Delta}^{(l)},
  \label{eq:emanation}
\end{equation}
where offsets are scaled by $1/\sqrt{l}$ for layer index~$l$.  The
Steersman monitors a \emph{coherence} metric---the mean Frobenius norm
of the offsets relative to the master---and applies regularization
pressure when individual layers fragment excessively from the global
topology.  The same RD co-axial pairs are specified, but contrastive
loss operates through the master matrix rather than per-layer bridges.
This tests whether global topological coherence can coexist with local
layer-wise adaptation.

\subsection{Spectral-Only Baseline}
\label{sec:spectral_baseline}

Channel counts $n \in \{3, 4, 6, 8, 12\}$ are trained with spectral
regularization alone---contrastive loss disabled.  The Steersman's
feedback loop targets Fiedler algebraic connectivity
$\lambda_2 \approx 0.1$ but provides no directional preference
among channel pairs.  This baseline establishes the \emph{spectral
attractor}: the bridge's unprogrammed equilibrium, the connectivity
level it reaches when the Steersman provides coupling pressure but
no topological specification.

\subsection{Training Setup}
\label{sec:training_setup}

All experiments use TinyLlama-1.1B-Chat with identical hyperparameters:
learning rate $2 \times 10^{-4}$, batch size~2, gradient
accumulation~8 (effective batch~16), LoRA rank~24, identity
initialization, random seed~42, maximum sequence length~512,
alpaca-cleaned dataset, 10{,}000~training steps.  The channel size
$s = r/n$ varies with~$n$: $s{=}8$ ($n{=}3$), $s{=}6$ ($n{=}4$),
$s{=}4$ ($n{=}6$), $s{=}3$ ($n{=}8$), $s{=}1$ ($n{=}24$).
Table~\ref{tab:topology_specs} summarizes all pair specifications.

\begin{table}[t]
\centering
\caption{Pair specifications for all experimental topologies.  Co-axial
  pairs receive positive contrastive gradient; all remaining
  off-diagonal pairs are cross-axial and receive negative gradient.
  $|\mathcal{C}|$ = co-axial pair count,
  $|\mathcal{X}|$ = cross-axial pair count.}
\label{tab:topology_specs}
\small
\begin{tabular}{llclcc}
\toprule
Topology & $n$ & Co-axial pairs & Source & $|\mathcal{C}|$ & $|\mathcal{X}|$ \\
\midrule
Octahedron (3D) & 4 & $(0,1), (2,3)$ & Vertex opposition & 2 & 4 \\
RD (3D)         & 6 & $(0,5), (1,4), (2,3)$ & Face opposition & 3 & 12 \\
Tesseract (4D)  & 8 & $(0,1), (2,3), (4,5), (6,7)$ & Vertex opposition & 4 & 24 \\
24-cell (4D)    & 24 & 12 antipodal pairs (Eq.~\ref{eq:24c_pairs}) & D4 root system & 12 & 264 \\
\midrule
Wrong-labels    & 6 & Random partition (seed 42) & None (control)  & 3 & 12 \\
Resonance       & 6 & $(0,2), (3,4), (0,5)$ & Number-theoretic & 3 & 12 \\
\bottomrule
\end{tabular}
\end{table}
